\section{Reductions and boundary conditions}\label{sec:reductions}

Three exact reductions are useful for interpreting the full calculation.

\subsection{Classical piezoelectric-semiconductor limit}
For \(\ell=H=0\), the strain-gradient and carrier--gradient sectors vanish and the
model reduces to the usual piezoelectric-semiconductor equations with mobile-carrier
attenuation \cite{hutson1962}. This limit is used as a numerical control; it is not
a claim that the flexoelectric coefficient is physically zero.

\subsection{Quasi-static electrostatic limit}
For \(\omega d/c_s\ll1\), where \(c_s=\sqrt{C_{44}/\rho}\), the electric potential
is eliminated algebraically at each Fourier mode. The model then retains
mechanical, strain-gradient, and carrier fields without adding an electromagnetic
radiation mode. The present calculation is in this quasistatic regime.

\subsection{Carrier-drift reduction}
For \(M=0\), the drift--diffusion operator becomes a diffusion operator and no
carrier current is generated. This is a control limit, not a physically calibrated
mobility choice.

\subsection{Variational boundary conditions}
On the plate faces \(x_3=\pm d/2\), the mechanical faces are traction-free,
\[
\sigma_{i3}=0,
\]
including the natural corner-line force associated with the strain-gradient term.
The electrical faces are ideally grounded in the primary case,
\[
\phi=0\quad(x_3=\pm d/2),
\]
and the corresponding open-circuit comparison uses
\[
D_3=0\quad(x_3=\pm d/2).
\]
The carrier faces are blocking in the primary case,
\[
J_3=0\quad(x_3=\pm d/2),
\]
which is imposed on the flux \(\mathbf J\), not by additionally prescribing
\(\Delta n=0\).

At the edge \(x=0\), the incident and outgoing angular-spectrum fields must satisfy
the same tangential \(k_y\), the mechanical traction/corner-force balance, the
electrical potential condition, and the carrier-flux condition. The edge solution
is obtained by matching all outgoing half-space modes, not by selecting a single
root. Backward classification requires three independent pieces of evidence: the
solved outgoing coefficient, negative thickness-integrated energy components, and
negative power along the physical backward direction. A mode is not classified as
backward solely from the sign of \(\Re K_x\).

\cref{tab:parameters} summarizes the primary parameters. The carrier pair
\(n_0=2\times10^{18}\,\)m\(^{-3}\) and \(M=0.0175\,\)m\(^2\)/(V\,s) is taken from a
primary bulk hydrothermal ZnO Hall measurement with stated transfer assumptions
\cite{suscavage1999}; the diffusivity is derived from the Einstein relation, not
measured. The elastic, piezoelectric, dielectric, and base carrier constants are
the retained literature-model values from Zhang and co-workers \cite{zhang2025}.
The gradient length and \(H\)-normalized values are synthetic, and the
flexoelectric coefficient is set to zero in the computed control.

\begin{table}[htbp]
\centering
\caption{Retained physical constants and provenance. The literature-model values
are not a joint same-specimen uncertainty analysis.}
\label{tab:parameters}
\begin{tabular}{ll}
\toprule
Constant & Value / provenance \\
\midrule
\(C_{11},C_{12},C_{13},C_{33},C_{44},C_{66}\) & 210, 121, 105, 221, 43, 44.5 GPa \\
\(e_{15},e_{31},e_{33}\) & \(-0.48\), \(-0.57\), \(1.32\) C/m\(^2\) \\
\(\kappa_{11},\kappa_{33}\) & \(7.61\times10^{-11}\), \(8.85\times10^{-11}\) F/m \\
\(\rho\) & \SI{5700}{\kilogram\per\cubic\metre} \\
Base \(n_0\), \(M\), \(D\) & \SI{2e18}{\per\cubic\metre}, \SI{0.0175}{\per\volt\per\second}, \SI{4.55e-4}{\square\metre\per\second} \\
\(\ell\), \(A\), \(H\) & synthetic; \(f=0\) control \\
\bottomrule
\end{tabular}
\end{table}
